\section{Problem formulation and notation}\label{sec:problem}
Let $P\succeq0$ be symmetric, $q$ finite, $A$ sparse and $x\in\R^n$ the
original decision. Each row value $z_i=A_i x$ has an interval
$[\ell_i,u_i]$ whose endpoints may be explicitly absent. Hard rows form
$\mathcal H$; all other admitted rows have a specified convex price. The
model is
\begin{equation}\label{eq:paper-model}
 \min_x\ \tfrac12x^\top P x+q^\top x+g(Ax)
 \quad\text{subject to}\quad
 \ell_i\leq A_ix\leq u_i\quad(i\in\mathcal H).
\end{equation}
The function $g$ is separable for the first method and a sum over disjoint
directional groups for the second. Both constructions keep physical hard
policy distinct from numerical penalty parameters. An elastic optimum can
purchase a preferred-row violation; that does not change what constitutes
hard feasibility.

Lower and upper departures are $v_{L,i}=(\ell_i-z_i)_+$ and
$v_{U,i}=(z_i-u_i)_+$. Corrective slack has sign
$s_i=v_{L,i}-v_{U,i}$, opposite to the signed violation residual
$z_i-\Pi_C(z)_i$. The distinction matters when conjugates and original
row multipliers are recovered. Upper finite-side multipliers enter the
original dual positively; lower finite-side multipliers enter negatively.
The fixed positive diagonal proximal metric is $H=\diag(h_i)$, $h_i>0$.
It is not a changing physical price or an arbitrary dense coupling matrix.

\begin{table}[htbp]\centering\small
\caption{Symbols and dimensions for the two proximal constructions. Physical row and decision units belong to each original model; computational duration and model prediction time remain different quantities.}\label{tab:paper-notation}
\begin{tabularx}{\textwidth}{@{}lXX@{}}\toprule
Symbol & Meaning & Dimension\\\midrule
$x,z=Ax$ & Decision and original row value & Decision / row unit\\
$C,\Pi_C,r$ & Interval box, box projection and residual & Row space / row unit\\
$a_i,b_i$ & Linear and quadratic side price & Cost/row unit; cost/row unit squared\\
$c_{U,i},c_{L,i}$ & Directional scale & Peak unit/row unit\\
$\Psi,g_C$ & Residual and eliminated interval penalty & Cost\\
$H,\alpha$ & Fixed proximal metric and relaxation & Cost/row unit squared; dimensionless\\
$y,\sigma_C,\Psi^*$ & Original multiplier, support and conjugate & Cost/row unit; cost; cost\\
$T,R$ & Computational duration and relative reduction & Seconds; dimensionless\\\bottomrule
\end{tabularx}
\end{table}

Table~\ref{tab:paper-notation} fixes the interpretation used in derivations
and figures. For a directional group, scales make physically different
departures comparable in one peak unit. With metres as row units and a
dimensionless peak, directional scales have inverse-metre units. A
reciprocal row conversion scales both endpoints and the scale coefficient;
changing only one directional coefficient instead changes the policy.
We do not add raw departures with incompatible dimensions or interpret a
population index as physical time.

Input validity does not imply solution existence. A hard model may be
infeasible; semidefinite curvature may admit unbounded directions or
nonunique optimizers. The splitting statements assume an attained saddle
point in addition to their closed-convex conditions. The native grouped
interface additionally requires nonempty boxes, disjoint groups and positive
finite directional scales. Mathematical extension to overlapping groups,
dense metrics or nonconvex prices is not inferred from tests of the
admitted cases. These assumptions keep the implemented objective, the
proximal proof and the experimental acceptance domain aligned.

The empirical quantities have their own dimensions: complete-call durations
are seconds, relative reductions are dimensionless and displayed as percent,
and counts identify observations or instances as specified by each protocol.
An unavailable result carries no numerical duration or objective for inference.
These conventions prevent missing values from becoming successful zero-cost
operations when tables or statistical arrays are assembled.
